% Forced split of Q4. Q4a (selected UV) is not invented.
% Q4b (minimal axial deformation forced by Thm. obs) is constructed
% and tree-level matched to Hehl-Datta. No Final Status [P] for Q4a.
\documentclass[11pt]{article}
\usepackage[margin=1.05in]{geometry}
\usepackage{amsmath,amssymb,amsthm}
\usepackage{enumitem,booktabs}
\usepackage[colorlinks=true,linkcolor=blue,citecolor=blue,urlcolor=blue]{hyperref}

\newtheorem{theorem}{Theorem}[section]
\newtheorem{proposition}[theorem]{Proposition}
\theoremstyle{definition}
\newtheorem{definition}[theorem]{Definition}
\newtheorem*{admission}{Admission}

\newcommand{\pP}{\textbf{[P]}}
\newcommand{\pC}{\textbf{[C]}}
\newcommand{\pO}{\textbf{[O]}}
\newcommand{\IB}{I_B}

\title{\textbf{A Minimal UV Deformation Forced by the Axial Obstruction\\
(Not a Selected Ultraviolet Completion)}}
\author{Robert W.\ McGwier\\
\small CTO, Cohere Technology Group\\
\small GitHub: \texttt{n4hy}}
\date{15 August 2026 --- Version 1}

\begin{document}
\maketitle

\begin{abstract}
\noindent
Paper~11 refused to invent a CFT, a string compactification, or an
asymptotically safe fixed point to fill the ultraviolet line of the
Emergent Gravity Program. That refusal stands. SuperGrok does not
repeal it.

Paper~14 nevertheless proved an obstruction: Einstein--Cartan
Hehl--Datta is a contact, while CFT axial Fisher is nonlocal. Any
local bulk that matches the latter \emph{and} reduces to the New
Standard Model in the infrared must therefore \emph{not} be
Einstein--Cartan at all scales. The unique quadratic deformation
that does this, with the M9.1 coefficient locked in the infrared,
is a massive axial field whose on-shell effective kernel is
\begin{equation*}
r(k)=\frac{3/16}{1+k^2/M^2}.
\end{equation*}
Tree-level matching is derived and machine-checked. The
position-space kernel is the four-dimensional Yukawa function, not
a delta. The massless limit kills the contact.

This is a deformation, not a selection. It does not produce the
Standard Model, it does not renormalize the metric, and it does
not name a holographic pair. Final Status ``UV completion'' remains
\pO{} for that selected-pair question. The axial deformation itself
is \pP{} as tree-level effective-field-theory algebra and \pC{} as
the physical identification of $S_\mu$ with axial torsion.

No \pO{} is moved to \pP{} for Q4a.
\end{abstract}

\section*{Standard of proof}

``A UV that works'' is split before any Lagrangian is written.

\begin{enumerate}[leftmargin=1.6em]
\item \textbf{Q4a (selected completion).} A microscopic theory whose
low-energy limit \emph{is} the New Standard Model, including the
installed gauge group, representations, Yukawas, and a
ultraviolet-complete metric sector. Inventing one to fill a
section is a fabrication \cite{McGwierXI}. Still \pO{}.
\item \textbf{Q4b (forced deformation).} The minimal local bulk
extension that (i)~reproduces
$\mathcal{L}_{\mathrm{HD}}=-(3\kappa/16)\,J_5\cdot J_5$ as
$M\to\infty$, and (ii)~has a nonlocal axial quadratic form at
finite $M$, so that Paper~14 Theorem~obs is not a contradiction.
This is an inverse problem with a unique \emph{quadratic}
solution once the infrared lock is imposed. That uniqueness is
inside the ansatz, not in the space of all quantum gravities.
\end{enumerate}

Numerical agreement is not proof. The algebra below is the proof
of Q4b at tree level. A script is a check.

\section{Why Einstein--Cartan cannot be its own UV}

The New Standard Model after torsion elimination is
\cite{HehlDatta1971,McGwierIII}
\begin{equation}
\mathcal{L}_{\mathrm{IR}}
=
\frac{1}{2\kappa}\mathring{R}
-\Lambda
+\mathcal{L}_{\mathrm{SM}}
-\frac{3\kappa}{16}\,J_5^\mu J_{5\mu},
\qquad
\kappa=8\pi G.
\label{eq:IR}
\end{equation}
The last term is ultralocal. In any unitary CFT in $d\ge 2$,
\begin{equation}
\langle J_5^\mu(x)J_5^\nu(y)\rangle
\sim
\frac{C_J}{\lvert x-y\rvert^{2d-2}}
\label{eq:JJ}
\end{equation}
is not. Paper~14 Theorem~obs is the statement that
\eqref{eq:IR} cannot equal ball relative entropy through second
order for generic axial sources. Dictionary improvement at linear
order \cite{McGwierVI} does not install a bulk kinetic term.

Therefore a holographic UV that keeps a CFT axial current must
contain a propagating bulk field dual to that current
\cite{Witten1998}. That field is not the Einstein--Cartan
connection. Einstein--Cartan is an infrared limit, or it is the
wrong theory for the axial channel.

\section{The unique quadratic deformation}

Let $S_\mu$ be an axial vector (the axial torsion dual, or an
independent Proca field: the algebra below does not care). The
most general local quadratic action with at most two derivatives,
a mass term, and a linear Dirac source is, in momentum space,
\begin{equation}
\mathcal{L}_k
=
A\Bigl(1+\frac{k^2}{M^2}\Bigr)\lvert S(k)\rvert^2
+
B\,S(-k)\cdot J_5(k).
\label{eq:UV}
\end{equation}
$A$ has the dimension of a stiffness; $M$ is the new mass;
$B$ is the Dirac--axial coupling. Completing the square,
\begin{equation}
S(k)
=
-\frac{B}{2A(1+k^2/M^2)}\,J_5(k),
\qquad
\mathcal{L}_{\mathrm{eff}}(k)
=
-\frac{B^2}{4A(1+k^2/M^2)}\,J_5\cdot J_5.
\label{eq:on}
\end{equation}
The infrared lock is the M9.1 theorem
\cite{HehlDatta1971}:
\begin{equation}
\frac{B^2}{4A}=\frac{3\kappa}{16}.
\label{eq:lock}
\end{equation}
Then, and only then,
\begin{equation}
r(k)
:=
\frac{\mathcal{L}_{\mathrm{eff}}(k)}{-\kappa\,J_5\cdot J_5}
=
\frac{3/16}{1+k^2/M^2}.
\label{eq:r}
\end{equation}
Without \eqref{eq:lock} the infrared coefficient is a free
parameter and M9.1 is lost. With it, the only remaining freedom
in this ansatz is the mass $M$.

\begin{theorem}[IR recovery]
\label{thm:ir}
Under \eqref{eq:UV} and \eqref{eq:lock},
$\lim_{M\to\infty}r(k)=3/16$ at every fixed $k$, and
$r(0)=3/16$ at every finite $M$. \pP{} (algebra).
\end{theorem}

\begin{theorem}[evasion of the contact obstruction]
\label{thm:nl}
At finite $M$, $r(k)$ is not constant. The position-space kernel
is the Euclidean four-dimensional Yukawa function
\begin{equation}
G(x)
=
\int\frac{d^4k}{(2\pi)^4}
\frac{e^{ik\cdot x}}{k^2+M^2}
=
\frac{M}{4\pi^2 r}\,K_1(Mr),
\qquad
r=\lvert x\rvert,
\label{eq:yuk}
\end{equation}
which satisfies $(-\square+M^2)G=\delta^{(4)}$ and is not
$\delta^{(4)}$ itself. \pP{} (Fourier analysis).
\end{theorem}

\begin{theorem}[massless degeneration]
\label{thm:m0}
For any fixed $k\neq 0$, $\lim_{M\to 0}r(k)=0$. A massless axial
gauge field does not produce the Hehl--Datta contact. \pP{}.
\end{theorem}

Theorem~\ref{thm:m0} is the holographic tension made sharp. A
\emph{conserved} CFT current is dual to a \emph{massless} bulk
gauge field \cite{Witten1998}. That theory is Einstein--Maxwell
(axial), not Einstein--Cartan, and it has no contact infrared.
A heavy field ($M\sim M_{\mathrm{Pl}}$) matches \eqref{eq:IR} in
the infrared and matches a \emph{non-conserved} or
high-dimension operator in the ultraviolet, not a conserved
current at every scale. The obstruction is evaded only by
changing the theory, exactly as Paper~14 said.

\section{What $M$ is allowed to be}

If $M$ is a few inverse fermis, spacetime torsion waves are
laboratory fields. That is a different Lagrangian and is already
excluded as Einstein--Cartan \cite{Hehl1976,Shapiro2002}. The
deformation is compatible with the program only if
\begin{equation}
M \gtrsim M_{\mathrm{Pl}},
\end{equation}
so that torsion propagation is a Planckian correction and
\eqref{eq:IR} is the effective theory below that scale. This is
a consistency requirement, not a derivation of $M$. Ghost-free
quadratic Poincar\'e-gauge kinetic terms exist in a restricted
parameter set \cite{Sezgin1980}; this note does not re-derive
that classification and does not claim that every quadratic
torsion kinetic term is healthy.

\section{Machine check}

Pre-registered gates C1--C5 were locked before the scripts ran.
Solver: \texttt{m9\_4\_uv\_axial.py}. Auditor (no solver import,
sympy completion plus Bessel identities):
\texttt{m9\_4\_audit\_uv.py}. Both report PASS on C1--C5.
The auditor confirms that without the infrared lock
\eqref{eq:lock} the coefficient is not $3/16$.

This is tree-level effective-field-theory algebra. It is not a
loop computation and it is not a holographic Ward-identity
calculation.

\section{What this does not do}

\begin{itemize}[leftmargin=1.4em]
\item It does not select $SU(3)_c\times SU(2)_L\times U(1)_Y$,
three generations, Yukawas, or $\Lambda$. Those remain
empirical input \cite{McGwierIII}.
\item It does not renormalize the metric sector. Einstein--Hilbert
is still non-renormalizable. Quadratic curvature / Stelle /
asymptotic safety are other Lagrangians, not derived here.
\item It does not name a CFT. The AdS dictionary is cited
\cite{Witten1998}, not constructed.
\item It does not repair de~Sitter (Paper~IX).
\item It does not make spacetime Hehl--Datta a laboratory field.
\item It does not move a \texttt{MODELS.md} cell.
\end{itemize}

\begin{admission}
I did not invent a UV completion of the New Standard Model. I
constructed the unique quadratic axial deformation that
preserves M9.1 and evades the contact obstruction, and I checked
the tree-level algebra two ways. Q4a remains \pO{}. Calling Q4b
``a UV that works'' is allowed only as infrared matching plus
nonlocality, never as a selected microscopic theory.
\end{admission}

\section{Verdict}

\begin{center}
\begin{tabular}{@{}p{0.46\textwidth}p{0.46\textwidth}@{}}
\toprule
Claim & Status \\
\midrule
Q4a selected pair / SM / metric UV & still \pO{}; not invented \\
Quadratic axial deformation \eqref{eq:UV}--\eqref{eq:r} & derived \\
Infrared lock recovers $3/16$ & \pP{}, audited \\
Yukawa kernel, not a contact & \pP{}, audited \\
Massless limit is not Hehl--Datta & \pP{} \\
Identification $S_\mu=$ axial torsion & \pC{} (physical reading) \\
Final Status ``UV completion'' as a selected theory & unchanged, \pO{} \\
\bottomrule
\end{tabular}
\end{center}

The phrase that is true: \emph{the infrared of the unique
quadratic axial deformation is the New Standard Model axial
sector, and the ultraviolet of that deformation is not
Einstein--Cartan.} The phrase that is false: \emph{we now have
the UV completion.}

\begin{thebibliography}{99}
\bibitem{HehlDatta1971}
F.W.~Hehl and B.K.~Datta,
J.\ Math.\ Phys.\ 12 (1971) 1334.
\bibitem{Hehl1976}
F.W.~Hehl, P.~von der Heyde, G.D.~Kerlick and J.M.~Nester,
Rev.\ Mod.\ Phys.\ 48 (1976) 393.
\bibitem{Sezgin1980}
E.~Sezgin and P.~van Nieuwenhuizen,
Phys.\ Rev.\ D 21 (1980) 3269.
\bibitem{Witten1998}
E.~Witten,
Adv.\ Theor.\ Math.\ Phys.\ 2 (1998) 253, arXiv:hep-th/9802150.
\bibitem{Shapiro2002}
I.L.~Shapiro,
Phys.\ Rept.\ 357 (2002) 113, arXiv:hep-th/0103093.
\bibitem{McGwierIII}
R.W.~McGwier,
``The New Standard Model for Unified Physics,''
\url{https://github.com/n4hy/New_Model_Emergent_Gravity}.
\bibitem{McGwierVI}
R.W.~McGwier,
``Evaluation of the Universal Kinematic Integral,'' ibid.
\bibitem{McGwierXI}
R.W.~McGwier,
``Open Questions, Attempts, and Admissions,'' ibid.
\bibitem{McGwierXIV}
R.W.~McGwier,
``Second-Order Einstein--Cartan from Entanglement,'' ibid.
\end{thebibliography}

\end{document}
